\documentclass[11pt,a4paper]{article}

\usepackage{fontspec}
\usepackage[english]{babel}
\babelprovide[import,onchar=ids fonts]{chinese-simplified}
\setmainfont{DejaVu Sans}
\setsansfont{DejaVu Sans}
\setmonofont{DejaVu Sans Mono}
\babelfont[chinese-simplified]{rm}[AutoFakeBold=2.0]{Droid Sans Fallback}
\babelfont[chinese-simplified]{sf}[AutoFakeBold=2.0]{Droid Sans Fallback}

\usepackage{amsmath}
\usepackage{amssymb}
\usepackage{mathtools}
\usepackage{geometry}
\usepackage{xcolor}
\usepackage{booktabs}
\usepackage{tabularx}
\usepackage{array}
\usepackage{enumitem}
\usepackage{listings}
\usepackage{fancyhdr}
\usepackage{hyperref}

\geometry{
  left=24mm,
  right=24mm,
  top=23mm,
  bottom=24mm
}

\definecolor{accent}{HTML}{185A78}
\definecolor{accentdark}{HTML}{123B50}
\definecolor{softblue}{HTML}{EAF3F7}
\definecolor{softgray}{HTML}{F3F5F6}
\definecolor{textgray}{HTML}{3F4A50}
\definecolor{codegreen}{HTML}{276749}
\definecolor{codepurple}{HTML}{6B46C1}

\hypersetup{
  colorlinks=true,
  linkcolor=accentdark,
  urlcolor=accent,
  citecolor=accentdark,
  pdfauthor={DiffusionNFT Project},
  pdftitle={Reward-Conditioned Branch-Free Adaptive Weight 设计说明}
}

\pagestyle{fancy}
\fancyhf{}
\fancyhead[L]{\small\color{textgray}DiffusionNFT Adaptive Weight}
\fancyhead[R]{\small\color{textgray}技术说明}
\fancyfoot[C]{\small\thepage}
\renewcommand{\headrulewidth}{0.4pt}
\renewcommand{\headrule}{\hbox to\headwidth{\color{accent}\leaders\hrule height \headrulewidth\hfill}}

\setlength{\parindent}{2em}
\setlength{\parskip}{0.35em}
\setlength{\headheight}{22pt}
\addtolength{\topmargin}{-7pt}
\setlength{\emergencystretch}{2em}
\setlist[itemize]{leftmargin=2em,itemsep=0.25em,topsep=0.3em}
\setlist[enumerate]{leftmargin=2.2em,itemsep=0.3em,topsep=0.3em}
\renewcommand{\contentsname}{目录}
\renewcommand{\tablename}{表}

\newcommand{\sg}{\operatorname{sg}}
\newcommand{\mabs}{\operatorname{m}}
\newcommand{\xcur}{\mathbf{x}_{0,i}^{\theta}}
\newcommand{\xold}{\mathbf{x}_{0,i}^{\mathrm{old}}}
\newcommand{\xclean}{\mathbf{x}_{0,i}}
\newcommand{\xtarget}{\mathbf{x}_{0,i}^{\mathrm{NFT}}}
\newcommand{\eps}{\epsilon}
\newcommand{\R}{\mathrm{RC}}

\newcolumntype{Y}{>{\raggedright\arraybackslash}X}

\lstdefinestyle{pythonstyle}{
  language=Python,
  basicstyle=\ttfamily\small,
  keywordstyle=\color{codepurple}\bfseries,
  commentstyle=\color{codegreen},
  stringstyle=\color{accentdark},
  backgroundcolor=\color{softgray},
  frame=single,
  rulecolor=\color{softgray},
  breaklines=true,
  breakatwhitespace=false,
  columns=fullflexible,
  keepspaces=true,
  showstringspaces=false,
  tabsize=4,
  xleftmargin=0.6em,
  xrightmargin=0.6em,
  aboveskip=0.8em,
  belowskip=0.8em
}

\begin{document}
\renewcommand{\contentsname}{目录}
\renewcommand{\tablename}{表}
\hypersetup{pageanchor=false}

\begin{titlepage}
\thispagestyle{empty}
\vspace*{28mm}
\noindent{\color{accent}\rule{\textwidth}{2.2pt}}
\vspace{10mm}

{\sffamily\bfseries\fontsize{23}{31}\selectfont
Reward-Conditioned Branch-Free\\[2mm]
Adaptive Weight 设计说明}

\vspace{8mm}
{\Large 面向 DiffusionNFT single-branch pseudo-target objective}

\vspace{17mm}
{\large
从原版双分支归一化出发，构造一个只需要单个 reward-conditioned
sample weight 的近似形式}

\vfill

\begin{tabular}{@{}ll@{}}
\textbf{对应推导} & \texttt{assets/ours\_idea.tex}\\
\textbf{设计版本} & Reward-conditioned root-second-moment normalization\\
\textbf{日期} & 2026 年 7 月 27 日
\end{tabular}

\vspace{15mm}
\noindent{\color{accent}\rule{\textwidth}{0.8pt}}
\end{titlepage}

\hypersetup{pageanchor=true}
\pagenumbering{roman}
\tableofcontents
\clearpage
\pagenumbering{arabic}

\section{核心结论}

新的 adaptive weight 不是把 single-branch loss 自身的误差直接放到分母，
而是根据 reward、old-to-clean 残差和 current-to-old 更新量，估计一个
\textbf{reward-conditioned residual second moment}。对第 $i$ 个样本和第
$j$ 个维度，定义
\begin{align}
d_{i,j}
&=
\left(\xold-\xclean\right)_j,
&
\Delta_{i,j}
&=
\left(\xcur-\xold\right)_j,
&
s_i
&=2r_i-1.
\end{align}
新的逐元素尺度为
\begin{equation}
\boxed{
q_{i,j}
=
\left(d_{i,j}+\beta s_i\Delta_{i,j}\right)^2
+
4\beta^2r_i(1-r_i)\Delta_{i,j}^2
}
\label{eq:q-main}
\end{equation}
并构造一个 detached 的 sample-level normalization factor：
\begin{equation}
\boxed{
w_i^{\R}(t)
=
\sg\!\left[
\max\!\left(
\frac{1}{D}\sum_{j=1}^{D}\sqrt{q_{i,j}},
\eps
\right)
\right].
}
\label{eq:w-main}
\end{equation}
最终的 single-branch objective 是
\begin{equation}
\boxed{
\ell_{i,\R}
=
\frac{\beta}{w_i^{\R}(t)}
\left\|
\xcur-\xtarget
\right\|^2.
}
\label{eq:loss-main}
\end{equation}

这里需要区分两个容易混淆的概念：
\begin{itemize}
  \item $w_i^{\R}$ 是放在分母中的 \textbf{normalization factor}；
  \item 真正乘在该样本 MSE 前面的 \textbf{effective sample weight} 是
  $\kappa_i^{\R}=\beta/w_i^{\R}$。
\end{itemize}

这个设计的核心价值是：它只需要一个 sample weight，不构造正负预测，
但仍保留原版 DiffusionNFT adaptive factor 最重要的残差归一化性质。
它在 $r_i=0$ 和 $r_i=1$ 两个端点与原版严格一致，在
$\xcur=\xold$ 的 policy refresh 点与原版具有严格相同的局部梯度；
对于中间 reward，它是一个有意设计的稳定近似。

\section{问题背景}

\subsection{符号与 pseudo-target}

DiffusionNFT 的 single-branch pseudo-target 可以写成
\begin{equation}
\xtarget
=
\xold
+
\frac{2r_i-1}{\beta}
\left(\xclean-\xold\right).
\label{eq:pseudo-target}
\end{equation}
使用
\begin{equation}
d_i=\xold-\xclean,
\qquad
\Delta_i=\xcur-\xold,
\qquad
s_i=2r_i-1,
\end{equation}
可将其改写为
\begin{equation}
\xtarget=\xold-\frac{s_i}{\beta}d_i,
\qquad
\xcur-\xtarget
=
\Delta_i+\frac{s_i}{\beta}d_i.
\label{eq:target-residual}
\end{equation}

$s_i\in[-1,1]$ 是 centered reward：
\begin{itemize}
  \item $s_i>0$：pseudo-target 从 old prediction 向 clean sample 移动；
  \item $s_i<0$：pseudo-target 向相反方向移动；
  \item $s_i=0$：pseudo-target 等于 old prediction。
\end{itemize}

本文统一使用
\begin{equation}
\|\mathbf{u}\|^2=\frac{1}{D}\sum_{j=1}^{D}u_j^2,
\qquad
\mabs(\mathbf{u})=\frac{1}{D}\sum_{j=1}^{D}|u_j|.
\end{equation}

\subsection{原版 adaptive weighting 做了什么}

用 $d_i$ 和 $\Delta_i$ 表示原版的两个残差：
\begin{equation}
e_i^+=d_i+\beta\Delta_i,
\qquad
e_i^-=d_i-\beta\Delta_i.
\end{equation}
对应的 detached normalization factors 是
\begin{equation}
w_i^+
=
\sg\!\left[\max\!\left(\mabs(e_i^+),\eps\right)\right],
\qquad
w_i^-
=
\sg\!\left[\max\!\left(\mabs(e_i^-),\eps\right)\right].
\end{equation}
原版 adaptive objective 为
\begin{equation}
\ell_{i,\mathrm{sep}}
=
\frac{1}{\beta}
\left[
\frac{r_i}{w_i^+}\|e_i^+\|^2
+
\frac{1-r_i}{w_i^-}\|e_i^-\|^2
\right].
\label{eq:original}
\end{equation}

这两个分母有两个作用。第一，它们将较大的重建误差除以较大的尺度，
避免困难样本或困难 timestep 完全主导梯度。第二，当 $w_i^+\neq w_i^-$
时，它们会改变两个方向的相对贡献，因此不仅改变梯度大小，还会改变严格
合并后 pseudo-target 的位置。

如果要求与式 \eqref{eq:original} 对任意中间 reward 严格等价，就必须保留
$w_i^+$、$w_i^-$ 的信息。严格合并后的梯度系数包含
\begin{equation}
\beta
\left(
\frac{r_i}{w_i^+}
+
\frac{1-r_i}{w_i^-}
\right),
\end{equation}
而 effective reward 还会变成
\begin{equation}
r_i^{\mathrm{eff}}
=
\frac{r_i/w_i^+}
{r_i/w_i^+ + (1-r_i)/w_i^-}.
\end{equation}
因此，严格等价版不可避免地较复杂。

\section{为什么直接 post-merge normalization 不够理想}

当前简易版直接使用 single-branch residual：
\begin{equation}
w_i^{\mathrm{post}}
=
\sg\!\left[
\max\!\left(
\mabs(\xcur-\xtarget),
\eps
\right)
\right],
\qquad
\ell_{i,\mathrm{post}}
=
\frac{\beta}{w_i^{\mathrm{post}}}
\|\xcur-\xtarget\|^2.
\label{eq:post}
\end{equation}
它实现简单，但分母测量的是“当前模型距离优化目标还有多远”，而原版分母
测量的是“reward 所作用的重建残差有多大”。这两者并不是同一个量。

\subsection{问题一：拟合 target 时分母必然触发下限}

当 $\xcur\rightarrow\xtarget$ 时，
\begin{equation}
\mabs(\xcur-\xtarget)\rightarrow 0.
\end{equation}
因此，未截断的分母必然趋近于零，实际实现只能依赖 $\eps$。虽然
$\eps$ 避免了数值除零，但分母会从 residual-dependent normalization
退化为由人工常数控制的区域。

尤其当 $r_i=1/2$ 时，$s_i=0$ 且 $\xtarget=\xold$。训练通常从 current
model 接近 old model 的位置开始，因此一个几乎中性的样本很容易立即落入
分母下限区域，即使 old prediction 与 clean sample 之间仍存在很大误差。

\subsection{问题二：在 old-policy 点附近会弱化 reward 强度}

在 $\xcur=\xold$，即 $\Delta_i=0$ 时，式 \eqref{eq:target-residual}
给出
\begin{equation}
w_i^{\mathrm{post}}
=
\frac{|s_i|}{\beta}\mabs(d_i)
\qquad
\text{（忽略 $\eps$）}.
\end{equation}
因为 $w_i^{\mathrm{post}}$ 被 detach，其梯度为
\begin{align}
\left.
\nabla_{\xcur}\ell_{i,\mathrm{post}}
\right|_{\Delta_i=0}
&=
\frac{2\beta}{D\,w_i^{\mathrm{post}}}
\left(\frac{s_i}{\beta}d_i\right)
\\
&=
\frac{2\beta\,\operatorname{sign}(s_i)}
{D\,\mabs(d_i)}
d_i,
\qquad s_i\neq 0.
\label{eq:post-gradient}
\end{align}
可以看到，$|s_i|$ 被分母抵消了。除 $s_i=0$ 的精确点外，一个略高于平均
的样本和一个 reward 极高的样本，可能获得近似相同的初始归一化梯度大小。
reward 主要只剩下正负号，而不是连续的置信强度。

这与原版 DiffusionNFT 在 policy refresh 点的行为不同。原版此时
$w_i^+=w_i^-=\mabs(d_i)$，梯度连续地正比于 $s_i$。新的设计首先要恢复
这一性质。

\section{新 weight 的推导}

\subsection{设计目标}

我们希望一个新的 normalization factor 同时满足：
\begin{enumerate}
  \item 每个样本只计算一个 detached scalar weight；
  \item weight 显式依赖 reward，但不生成正负 branch prediction；
  \item 在 current model 等于 old model 时，恢复原版的局部梯度；
  \item 在 reward 两个端点恢复原版目标；
  \item 对中间 reward，不因 signed residual 的抵消而异常变小；
  \item 保留原版使用 mean absolute residual 进行尺度归一化的习惯。
\end{enumerate}

\subsection{从方向残差的二阶矩出发}

为了说明新公式与原版的关系，可以引入一个只用于推导的方向变量
$Z_i\in\{-1,+1\}$：
\begin{equation}
\Pr(Z_i=+1)=r_i,
\qquad
\Pr(Z_i=-1)=1-r_i.
\end{equation}
它的均值和方差为
\begin{equation}
\mathbb{E}[Z_i]=2r_i-1=s_i,
\qquad
\operatorname{Var}(Z_i)=1-s_i^2=4r_i(1-r_i).
\end{equation}
考虑第 $j$ 个维度上的方向残差
\begin{equation}
E_{i,j}(Z_i)=d_{i,j}+\beta Z_i\Delta_{i,j}.
\end{equation}
它的均值是
\begin{equation}
\mathbb{E}[E_{i,j}]
=
d_{i,j}+\beta s_i\Delta_{i,j},
\end{equation}
方差是
\begin{equation}
\operatorname{Var}(E_{i,j})
=
\beta^2\operatorname{Var}(Z_i)\Delta_{i,j}^2
=
4\beta^2r_i(1-r_i)\Delta_{i,j}^2.
\end{equation}
由“二阶矩等于均值平方加方差”得到
\begin{align}
\mathbb{E}\!\left[E_{i,j}^2\right]
&=
\left(
d_{i,j}+\beta s_i\Delta_{i,j}
\right)^2
+
4\beta^2r_i(1-r_i)\Delta_{i,j}^2
\\
&=q_{i,j}.
\end{align}

这就是式 \eqref{eq:q-main}。推导借助了原版的方向分布来解释它为什么合理，
但实际计算只使用 $d_i$、$\Delta_i$ 和 $r_i$ 三个量，不构造任何 branch
prediction，也不维护 branch-specific weight。

\subsection{为什么使用
\texorpdfstring{$\frac{1}{D}\sum_j\sqrt{q_{i,j}}$}{mean root-second-moment}}

$\sqrt{q_{i,j}}$ 是第 $j$ 个维度 residual 的 root second moment。
先逐元素开方再取平均，是为了尽量保留原版 mean absolute error 的尺度：
\begin{itemize}
  \item 当 reward 位于端点时，$\sqrt{q_{i,j}}$ 会严格退化为对应 residual
  的绝对值；
  \item 因此 $\frac{1}{D}\sum_j\sqrt{q_{i,j}}$ 会严格退化为原版的
  $\mabs(e_i^\pm)$；
  \item 如果改为 $\sqrt{\frac{1}{D}\sum_j q_{i,j}}$，得到的是 global RMS，
  在 reward 端点也无法与原版 MAE factor 严格一致。
\end{itemize}

\subsection{为什么必须 stop-gradient}

$w_i^{\R}$ 的目标是度量当前样本的更新尺度，而不是成为模型可以主动优化的
对象。因此必须使用 $\sg[\cdot]$ 或在代码中置于 \texttt{torch.no\_grad()}
上下文中。

如果不 detach，梯度会穿过分母。模型可能通过改变 $q_{i,j}$ 来改变自己的
loss weight，而不只是拟合 pseudo-target。这样不仅破坏与原版 adaptive
factor 的对应关系，还可能诱导模型增大某些 residual 以减小 effective
sample weight。

\section{三个关键性质}

\subsection{性质一：两个 reward 端点与原版严格一致}

当 $r_i=1$ 时，$s_i=1$，方差项为零：
\begin{equation}
q_{i,j}
=
\left(d_{i,j}+\beta\Delta_{i,j}\right)^2,
\qquad
w_i^{\R}=w_i^+.
\end{equation}
同时
\begin{equation}
\xcur-\xtarget
=
\Delta_i+\frac{1}{\beta}d_i.
\end{equation}
所以
\begin{equation}
\ell_{i,\R}
=
\frac{\beta}{w_i^+}
\left\|
\Delta_i+\frac{d_i}{\beta}
\right\|^2
=
\frac{1}{\beta w_i^+}
\|d_i+\beta\Delta_i\|^2,
\end{equation}
与式 \eqref{eq:original} 在 $r_i=1$ 时完全相同。

同理，当 $r_i=0$ 时，
\begin{equation}
q_{i,j}
=
\left(d_{i,j}-\beta\Delta_{i,j}\right)^2,
\qquad
w_i^{\R}=w_i^-,
\end{equation}
并且 $\ell_{i,\R}$ 与原版 $r_i=0$ 的目标完全相同。

\subsection{性质二：在 policy refresh 点与原版梯度严格一致}

当 current model 与 old model 相同时，$\Delta_i=0$。此时对任意 reward，
\begin{equation}
q_{i,j}=d_{i,j}^2,
\qquad
w_i^{\R}
=
\sg\!\left[\max\!\left(\mabs(d_i),\eps\right)\right]
=w_i^+=w_i^-.
\label{eq:refresh-weight}
\end{equation}

记公共 factor 为 $w_i^0$。由式 \eqref{eq:original} 得到原版梯度
\begin{align}
\left.
\nabla_{\xcur}\ell_{i,\mathrm{sep}}
\right|_{\Delta_i=0}
&=
\frac{2}{D}
\left[
\frac{r_i}{w_i^0}d_i
-
\frac{1-r_i}{w_i^0}d_i
\right]
\\
&=
\frac{2s_i}{D\,w_i^0}d_i.
\label{eq:original-local-grad}
\end{align}
新 objective 的梯度为
\begin{align}
\left.
\nabla_{\xcur}\ell_{i,\R}
\right|_{\Delta_i=0}
&=
\frac{2\beta}{D\,w_i^0}
\left(
\frac{s_i}{\beta}d_i
\right)
\\
&=
\frac{2s_i}{D\,w_i^0}d_i.
\label{eq:new-local-grad}
\end{align}
式 \eqref{eq:original-local-grad} 与式 \eqref{eq:new-local-grad} 完全相同。

这条性质很重要，因为每次 old policy 更新后，current 与 old 通常接近。
新 weight 在训练最常访问的局部区域保持了原版 adaptive objective 的
reward 响应：梯度大小连续地正比于 $s_i=2r_i-1$。它不会像 post-merge
版本那样把 $|s_i|$ 从初始梯度中归一化掉。

\subsection{性质三：中间 reward 拟合 target 时 factor 不会普遍塌缩}

当 current prediction 正好等于 pseudo-target 时，
\begin{equation}
\Delta_i^\star
=
-\frac{s_i}{\beta}d_i.
\end{equation}
代入式 \eqref{eq:q-main}：
\begin{align}
q_{i,j}^\star
&=
\left(1-s_i^2\right)^2d_{i,j}^2
+
\left(1-s_i^2\right)s_i^2d_{i,j}^2
\\
&=
\left(1-s_i^2\right)d_{i,j}^2.
\end{align}
因此在忽略 $\eps$ 时，
\begin{equation}
\boxed{
w_i^{\R}(\xcur=\xtarget)
=
\sqrt{1-s_i^2}\,\mabs(d_i)
=
2\sqrt{r_i(1-r_i)}\,\mabs(d_i).
}
\label{eq:target-scale}
\end{equation}

对任何严格的中间 reward $0<r_i<1$，只要 old-to-clean error 不为零，
这个 factor 就仍然为正。几个典型情况如下。

\begin{table}[ht]
\centering
\small
\caption{拟合 pseudo-target 时的新 normalization factor，表中省略
$\mabs(d_i)$ 和 $\eps$。}
\begin{tabular}{cccccc}
\toprule
$r_i$ & $0$ & $0.25$ & $0.5$ & $0.75$ & $1$\\
\midrule
$s_i=2r_i-1$ & $-1$ & $-0.5$ & $0$ & $0.5$ & $1$\\
$2\sqrt{r_i(1-r_i)}$ & $0$ & $0.866$ & $1$ & $0.866$ & $0$\\
\bottomrule
\end{tabular}
\end{table}

特别地，$r_i=0.5$ 时
\begin{equation}
w_i^{\R}=\mabs(d_i),
\end{equation}
而 post-merge factor 为零并触发 $\eps$。端点处新 factor 也会趋于零，
但这是合理的：此时唯一被 reward 选择的 residual 已经被完全拟合，原版
factor 同样会触发 $\eps$。

\section{如何理解 reward 的作用}

\subsection{均值项与方差项}

式 \eqref{eq:q-main} 可以分成两部分：
\begin{equation}
\underbrace{
\left(d_{i,j}+\beta s_i\Delta_{i,j}\right)^2
}_{\text{reward-directed mean residual}}
\quad+\quad
\underbrace{
4\beta^2r_i(1-r_i)\Delta_{i,j}^2
}_{\text{directional variance}}.
\end{equation}

均值项描述 reward 指定的净方向；方差项描述当 reward 不确定时，不能被
一个 signed mean residual 表示的剩余尺度。
\begin{itemize}
  \item 当 $r_i$ 接近 $0$ 或 $1$，reward 方向明确，方差项趋于零；
  \item 当 $r_i=0.5$，方向最不确定，方差系数达到最大值；
  \item 方差项始终非负，所以即使均值项发生向量抵消，也不会轻易得到一个
  不合理的小 factor。
\end{itemize}

\subsection{为什么 refresh 点的 factor 暂时不依赖 reward}

在 $\Delta_i=0$ 时，式 \eqref{eq:refresh-weight} 对所有 reward 都等于
$\mabs(d_i)$。这不是缺少 reward conditioning，而是有意保留原版的局部
几何结构。

此时 reward 已经通过 pseudo-target 控制梯度：
\begin{equation}
\nabla_{\xcur}\ell_{i,\R}
\propto
s_i d_i.
\end{equation}
normalization factor 只负责消除不同样本 $d_i$ 尺度的差异，不应再把
$|s_i|$ 消掉。当 current model 离开 old model 后，$\Delta_i\neq 0$，
$w_i^{\R}$ 会通过均值项和方差项显式依赖 reward。

\subsection{直接使用 normalized advantage 的形式}

如果后续方法使用归一化 reward 或 advantage
$\widehat{A}_i\in[-1,1]$，可以直接令
\begin{equation}
\widehat{A}_i=s_i=2r_i-1.
\end{equation}
此时新 factor 可以写成更适合 single-branch 方法的形式：
\begin{equation}
\boxed{
q_{i,j}
=
\left(
d_{i,j}+\beta\widehat{A}_i\Delta_{i,j}
\right)^2
+
\beta^2
\left(
1-\widehat{A}_i^2
\right)
\Delta_{i,j}^2.
}
\label{eq:advantage-form}
\end{equation}
这个表达式完全不需要正负 branch 的语义，只需要：
\begin{itemize}
  \item old prediction error $d_i$；
  \item current policy displacement $\Delta_i$；
  \item 每个样本的 normalized reward $\widehat{A}_i$。
\end{itemize}

如果使用论文中的线性 sample weight
$W_i=1+\widehat{A}_i\in[0,2]$，也可以直接使用
$\widehat{A}_i=W_i-1$ 代入式 \eqref{eq:advantage-form}。因此，这个
adaptive factor 可以自然接到后续“每个样本只有一个 reward weight”的
推导中。

\section{与另外两个版本的对比}

\begin{table}[ht]
\centering
\footnotesize
\caption{三种 adaptive weighting 形式的主要差异。}
\begin{tabularx}{\textwidth}{@{}p{0.19\textwidth}YYY@{}}
\toprule
 & 严格等价版 & 直接 post-merge 简易版 & 新 reward-conditioned 版\\
\midrule
Normalization factors
& 两个：$w_i^+,w_i^-$
& 一个：$w_i^{\mathrm{post}}$
& 一个：$w_i^{\R}$\\
\addlinespace
是否构造方向预测
& 是
& 否
& 否\\
\addlinespace
中间 reward 严格等价
& 是
& 否
& 否\\
\addlinespace
是否保留 adaptive target shift
& 是
& 否
& 否\\
\addlinespace
refresh 点局部梯度
& 原版基准
& 一般只保留 reward 符号
& 与原版严格一致\\
\addlinespace
拟合 pseudo-target 时的分母
& 由两个 residual 决定
& 普遍触发 $\eps$
& 中间 reward 通常保持非零\\
\addlinespace
reward 如何进入 factor
& 通过两个 factor 的混合
& 通过 pseudo-target residual
& 直接进入均值与方差\\
\addlinespace
计算与实现复杂度
& 最高
& 最低
& 介于两者之间，但仍是单 weight\\
\bottomrule
\end{tabularx}
\end{table}

严格等价版适合作为理论基准。post-merge 版最简单，但它会把 pseudo-target
距离和 adaptive normalization 混为同一个量。新版本的定位是：在不恢复
branch-specific target shift 的前提下，尽可能保留原版的 residual-scale
normalization、reward 连续性和关键边界行为。

\section{梯度与优化含义}

由于 $w_i^{\R}$ 被 detach，新 loss 对 current prediction 的梯度非常清晰：
\begin{equation}
\nabla_{\xcur}\ell_{i,\R}
=
\frac{2\beta}{D\,w_i^{\R}}
\left(
\xcur-\xtarget
\right).
\label{eq:new-gradient}
\end{equation}
因此：
\begin{itemize}
  \item \textbf{梯度方向}完全由 single-branch pseudo-target 决定；
  \item \textbf{梯度大小}由单个 sample multiplier
  $\kappa_i^{\R}=\beta/w_i^{\R}$ 调整；
  \item adaptive factor 不会像严格等价版那样再次移动 pseudo-target；
  \item 这一职责分离更符合后续方法的概念结构：归一化 reward 决定样本
  目标或样本权重；adaptive factor 只负责控制该样本的更新尺度。
\end{itemize}

对于 residual 很大的样本，$w_i^{\R}$ 通常也较大，因此更新被压缩；对于
residual 较小但未进入 $\eps$ 区域的样本，更新相对增强。这延续了原版
DiffusionNFT 的“按样本和 timestep 归一化优化难度”的思想。

\section{PyTorch 实现}

下面给出与公式直接对应的实现框架。关键要求是：用于 factor 的
\texttt{x0\_prediction} 必须在计算 $q_i$ 时 detach，但用于 MSE 的
\texttt{x0\_prediction} 必须保留梯度。

\begin{lstlisting}[style=pythonstyle]
# x0_prediction: current model clean prediction, with gradients
# x0_old_prediction, x0, reward: detached / no-grad tensors
x0_prediction = xt - t_expanded * forward_prediction
x0_old_prediction = xt - t_expanded * old_prediction.detach()

reduce_dims = tuple(range(1, x0.ndim))
r_view = r.view(-1, *([1] * (x0.ndim - 1)))

with torch.no_grad():
    d = (x0_old_prediction - x0).double()
    delta = (x0_prediction - x0_old_prediction).double()
    r64 = r_view.double()
    s = 2.0 * r64 - 1.0

    q = (
        (d + config.beta * s * delta).square()
        + 4.0
        * (config.beta ** 2)
        * r64
        * (1.0 - r64)
        * delta.square()
    )

    # Guard against a tiny negative value caused by finite precision.
    weight_factor = (
        q.clamp_min(0.0)
        .sqrt()
        .mean(dim=reduce_dims, keepdim=True)
        .clamp_min(1e-5)
        .to(x0_prediction.dtype)
    )

    x0_target = (
        x0_old_prediction
        + (s / config.beta).to(x0.dtype)
        * (x0 - x0_old_prediction)
    )

sample_loss = (
    config.beta
    * (x0_prediction - x0_target).square()
    / weight_factor
).mean(dim=reduce_dims)

policy_loss = sample_loss.mean()
\end{lstlisting}

\subsection{实现注意事项}

\begin{enumerate}
  \item \textbf{Reward 范围。}
  公式要求 $r_i\in[0,1]$，等价地要求 $s_i\in[-1,1]$。如果输入是未经
  限制的 advantage，应先标准化和裁剪。否则
  $4r_i(1-r_i)$ 或 $1-s_i^2$ 可能为负。

  \item \textbf{精度。}
  建议像原版 factor 一样使用 \texttt{float64} 计算 $q_i$ 和均值，再转回
  模型 dtype。特别是在 mixed precision 下，这能减少平方和开方带来的
  舍入问题。

  \item \textbf{非负保护。}
  理论上 $q_{i,j}\geq0$，因为它是两个平方项之和。代码中的
  \texttt{clamp\_min(0)} 只用于防止有限精度产生极小负值。

  \item \textbf{下限 $\eps$。}
  应与原版使用同一个 $\eps$，例如 $10^{-5}$，这样端点行为和 loss scale
  更可比。若改变 $\eps$，需同时观察 floor activation ratio。

  \item \textbf{按 sample 归约。}
  $w_i^{\R}$ 应对通道和空间维度取平均，但保留 batch 维。不要先对整个
  batch 求平均，否则会失去 sample-wise normalization。

  \item \textbf{日志。}
  建议记录 $w_i^{\R}$、$\beta/w_i^{\R}$、触发 $\eps$ 的样本比例，以及
  按 reward 和 timestep 分桶后的均值。
\end{enumerate}

\section{局限性与实验预期}

\subsection{它不在哪些方面等价}

对 $0<r_i<1$，新版本并不严格等价于原版。原因是原版的两个 detached
factor 会同时产生：
\begin{itemize}
  \item 一个由 $r_i/w_i^+ + (1-r_i)/w_i^-$ 决定的 harmonic-like
  gradient scale；
  \item 一个由 $w_i^+$ 与 $w_i^-$ 的差异引起的 adaptive pseudo-target
  shift。
\end{itemize}

单个 second-moment factor 无法同时保留这两个自由度。新设计主动放弃
adaptive target shift，只保留原始 single-branch pseudo-target，并用一个
reward-conditioned scalar 控制更新大小。

\subsection{为什么预期优于直接简易版}

预期改善来自以下三点，但最终仍需实验验证：
\begin{enumerate}
  \item 在 policy refresh 点恢复与原版相同的 reward-linear gradient，
  不会把 reward magnitude 归一化成近似符号；
  \item 对中间 reward，即使 current prediction 拟合 pseudo-target，
  factor 仍保留 old-to-clean residual 的尺度；
  \item 方差项阻止 signed mean residual 的偶然抵消，使 sample weight
  随 reward 和 policy displacement 更平滑地变化。
\end{enumerate}

新 factor 使用 root second moment。由 RMS--AM 关系，它可能比对应的
reward-weighted mean absolute residual 更保守，因此某些区间的
$\beta/w_i^{\R}$ 可能偏小。这不是数值错误，而是该近似用稳定性换取简洁性
的结果。若整体更新显著变小，应优先使用一个全局 loss-scale coefficient
统一校准，而不是破坏 $q_i$ 的 reward 结构。

\subsection{推荐消融}

为了判断它是否确实优于 post-merge 简易版，建议在完全相同的 rollout、
reward、$\beta$、学习率和随机种子下比较：
\begin{itemize}
  \item exact adaptive single-branch；
  \item direct post-merge adaptive；
  \item reward-conditioned branch-free adaptive；
  \item 不使用 adaptive factor 的 single-branch baseline。
\end{itemize}

除最终 reward 外，建议同时观察：
\begin{itemize}
  \item 训练前期不同 reward bin 的平均梯度范数；
  \item $w_i$ 与 effective multiplier $\beta/w_i$ 的分布；
  \item 分母触发 $\eps$ 的比例；
  \item 不同 timestep 的归一化尺度；
  \item current model 与 old model 的距离，以及训练是否出现突发梯度。
\end{itemize}

\section{总结}

新 adaptive weight 可以概括为以下思路：
\begin{quote}
\color{accentdark}
不再分别问“正方向 residual 多大”和“负方向 residual 多大”，而是根据
当前样本的 centered reward，直接估计 reward-directed residual 的
二阶尺度，并将它作为该样本唯一的 detached normalization factor。
\end{quote}

它比严格等价版更简单，因为只计算一个 scalar factor 且不改变
single-branch pseudo-target；它又比直接 post-merge 版更接近原版，因为
它恢复了 reward 端点、policy refresh 点和 residual-scale normalization
的关键行为。

最简洁的 normalized-advantage 表达为
\begin{equation}
\boxed{
w_i^{\R}
=
\sg\!\left[
\max\!\left(
\frac{1}{D}
\sum_{j=1}^{D}
\sqrt{
\left(d_{i,j}+\beta\widehat{A}_i\Delta_{i,j}\right)^2
+
\beta^2(1-\widehat{A}_i^2)\Delta_{i,j}^2
},
\eps
\right)
\right],
}
\end{equation}
其中 $\widehat{A}_i\in[-1,1]$。最终每个样本只有一个 effective weight
$\beta/w_i^{\R}$，这与后续基于 reward 设计 sample-wise weight 的方法
具有更一致的概念结构。

\end{document}
